\section{Geodesic passages and the index form}\label{sec:variations}

We now study the index form along a complete unit-speed geodesic of the
constructed metric.  The metric outside the neck has nonpositive sectional
curvature.  Inside the neck, concavity of the inverse metric gives a nonnegative
coordinate expression for the second variation.  The distinction between
that expression and the intrinsic index form is an endpoint term.  We first
compute that term and dispose of the passages that have no turn; the next
section estimates the remaining turning passages.

Throughout this section the neck metric is
\[
 g=dt^2+dx^{\mathsf T}Q(t)\,dx,\qquad P(t)=Q(t)^{-1},\qquad 0\le t\le1,
\]
with the smooth positive definite matrices constructed above.  In particular,
$-P''(t)$ is positive semidefinite.  At either boundary, we may use its inward
coordinate $\tau$ and its fixed torus frame.  In these coordinates,
\begin{equation}\label{eq:variational-end-germ}
 P(\tau)=\operatorname{diag}(1+\delta A(\tau),b)
 \quad(0\le\tau\le1/4),\qquad
 b=1+\delta A(1),\qquad \kappa=\delta A'(0)>0.
\end{equation}
These are exactly the two frames of the construction; changing between them
is a constant linear coordinate change, together with $\tau=1-t$ at the
right boundary.  The metric and these coordinates extend smoothly through
each boundary into a product collar.  The two end portions of the neck have
nonpositive sectional curvature.

\subsection{Normal index forms and coordinate second variation}

Let $\gamma$ be a complete unit-speed geodesic, and write $T=\dot\gamma$.
For a smooth field $V$ along $\gamma$ and an interval $J$, define
\begin{align}
 I_J(V)&=\int_J\bigl(|D_sV|^2-
       \langle R(V,T)T,V\rangle\bigr)\,ds,\label{eq:index-definition}\\
 q_V&=\langle V,T\rangle,\qquad
 I_J^\perp(V)=I_J(V)-\int_J|\dot q_V|^2\,ds.
       \label{eq:normal-index-definition}
\end{align}
We use the curvature convention for which the curvature term
$\langle R(V,T)T,V\rangle$ is nonpositive when sectional curvature is
nonpositive.  When $J$ is unbounded we will only use fields vanishing for
sufficiently large $|s|$, so these integrals are finite.

\begin{lemma}[Tangential invariance]\label{lem:tangential-invariance}
The integrand defining $I_J^\perp(V)$ is the index integrand of
$V^\perp=V-q_VT$.  Consequently, for every smooth real function $h$,
\[
 I_J^\perp(V+hT)=I_J^\perp(V)
\]
provided both fields have the finite integrals required in
\eqref{eq:normal-index-definition}.  This invariance holds pointwise, without
any condition on $h$ at finite endpoints.
\end{lemma}

\begin{proof}
Since $D_sT=0$ and $\langle V^\perp,T\rangle=0$, we have
$\langle D_sV^\perp,T\rangle=0$ and
\[
 |D_sV|^2=|D_sV^\perp|^2+|\dot q_V|^2.
\]
The curvature symmetries give
$\langle R(V,T)T,V\rangle=
\langle R(V^\perp,T)T,V^\perp\rangle$.
Both assertions follow.
\end{proof}

On any connected interval where $\gamma$ lies in the neck, lift its torus
coordinate to $\mathbb R^2$.  With a dot denoting differentiation in $s$ and
a prime differentiation in $t$, the geodesic equations imply
\begin{equation}\label{eq:neck-geodesic}
 p=Q(t)\dot x\ \text{is constant},\qquad
 F(t)=p^{\mathsf T}P(t)p,\qquad
 v=\dot t,\qquad
 v^2=1-F(t),\qquad \dot v=-\tfrac12F'(t).
\end{equation}
Indeed, the Euler--Lagrange equation in the cyclic variable $x$ gives the
first identity.  Unit speed gives the fourth, and the $t$ equation is
$\ddot t=\dot x^{\mathsf T}Q'\dot x/2=-F'/2$.
The function $F$ is concave because $P''\preceq0$.

For a field $V=(w,z)$ in these coordinates, $w$ is its interval component
and $z$ its two torus components.  Write $|\xi|_Q^2=\xi^{\mathsf T}Q\xi$.
On a compact coordinate interval $[a,c]$, define
\begin{equation}\label{eq:linear-energy}
 E_{[a,c]}(V)=\int_a^c
 \left(\dot w^2+|\dot z-P'p\,w|_Q^2
                 -\tfrac12F''(t)w^2\right)\,ds.
\end{equation}
Every term in this expression is nonnegative. We use the same integral
to define $E_J(V)$ on an unbounded neck interval $J$ when the field
vanishes for sufficiently large $|s|$.

\begin{lemma}[Endpoint correction]\label{lem:coordinate-index}
For a smooth field $V=(w,z)$ on $[a,c]$,
\begin{equation}\label{eq:coordinate-index}
 I_{[a,c]}(V)=E_{[a,c]}(V)
       -\bigl[\langle T,\Gamma(V,V)\rangle\bigr]_a^c,
\end{equation}
where $\Gamma$ denotes the Christoffel symbols in the chosen product
coordinates and brackets mean final value minus initial value.
In the inward coordinates at either neck boundary, if $w=0$ at that
boundary, then
\begin{equation}\label{eq:boundary-christoffel}
 \Gamma(V,V)=\tfrac\kappa2 z_1^2\partial_\tau.
\end{equation}
Thus the endpoint contribution in \eqref{eq:coordinate-index} is
nonnegative at an entry into or exit from the open neck, provided the
interval component of the field vanishes there.
\end{lemma}

\begin{proof}
Consider the coordinate-linear variation
$(t(s)+r w(s),x(s)+r z(s))$.  Its energy is one half the integral of its
squared speed.  Differentiating twice at $r=0$ gives
\[
 \int_a^c\left(\dot w^2+\dot z^{\mathsf T}Q\dot z
   +2w\dot z^{\mathsf T}Q'\dot x
   +\tfrac12w^2\dot x^{\mathsf T}Q''\dot x\right)\,ds.
\]
The inverse-matrix identities
\[
 Q'=-QP'Q,\qquad
 Q''=2QP'QP'Q-QP''Q,
\]
together with $\dot x=Pp$, turn this expression into
$E_{[a,c]}(V)$.

For completeness, the endpoint term can be seen directly in the
intrinsic second variation.  If $W=\partial_r\alpha$ and
$U=\partial_s\alpha$ for a variation $\alpha$, the first energy derivative
is $[\langle W,U\rangle]_a^c-\int\langle W,D_sU\rangle$.
Differentiate once more at the geodesic $r=0$, commute the covariant
derivatives, and integrate the $D_s^2V$ term by parts.  The result is
\[
 \frac{d^2}{dr^2}\bigg|_{r=0}\frac12\int_a^c|U|^2\,ds
 =I_{[a,c]}(V)+[\langle T,D_rW\rangle]_a^c.
\]
For the coordinate-linear variation, its coordinate second derivative in
$r$ is zero, so $D_rW=\Gamma(V,V)$.  This proves
\eqref{eq:coordinate-index}.

The nonzero types of Christoffel symbols for this product-coordinate
metric are
\[
 \Gamma^t_{ij}=-\tfrac12Q'_{ij},\qquad
 \Gamma^i_{tj}=\Gamma^i_{jt}=\tfrac12(PQ')^i{}_j.
\]
At a boundary where $w=0$, only
$-z^{\mathsf T}Q'z\,\partial_\tau/2$ remains.  By
\eqref{eq:variational-end-germ}, $Q'(0)=-\kappa e_1e_1^{\mathsf T}$;
hence \eqref{eq:boundary-christoffel} follows.  At an initial boundary
$\dot\tau>0$ and this scalar is added in
\eqref{eq:coordinate-index}; at a final boundary $\dot\tau<0$ and it is
subtracted.  Both signs are nonnegative.
\end{proof}

\subsection{The possible passages}

A \emph{neck passage} is a connected component $J$ of
$\{s\in\mathbb R:0<t(\gamma(s))<1\}$.  The notation means the preimage
of the open neck; it does not require an interval coordinate away from
the neck.  A restricted compactly supported field may have arbitrary
values at finite endpoints of $J$.

\begin{lemma}[Passage classification]\label{lem:neck-passages}
Every finite endpoint of a neck passage is transverse to the boundary
torus.  Exactly one of the following alternatives holds.
\begin{enumerate}
\item The passage is the whole real line and $t$ is constant.
\item The interval component $v=\dot t$ never vanishes.  The passage is
either a finite crossing between the two different boundary tori, or a
half-infinite interval approaching an interior equilibrium height.
\item After choosing the coordinates from one boundary and translating
the geodesic parameter, the passage has closure $[0,L]$, begins and ends
at that boundary, and has a unique turn at $s=L/2$.  In these coordinates
\begin{gather}
 t(0)=t(L)=0,\qquad
 v(0)=-v(L)=v_0>0,\qquad t_*=t(L/2)\in(0,1),\label{eq:turn-data}\\
 t(L-s)=t(s),\qquad v(L-s)=-v(s),\qquad
 0<F'(t)\le\kappa p_1^2\quad(0\le t\le t_*),\label{eq:turn-concavity}\\
 p_1\ne0,\qquad 0<v_0^2=F(t_*)-F(0)\le\kappa p_1^2.
 \label{eq:turn-incidence}
\end{gather}
\end{enumerate}
An interior equilibrium height $t_\infty$ in the second alternative
satisfies $F(t_\infty)=1$ and $F'(t_\infty)=0$; the velocity may tend to
zero there.  A geodesic tangent to a boundary torus at a finite time does
not enter the open neck at that time.
\end{lemma}

\begin{proof}
Use inward coordinates at a finite endpoint.  If $\dot\tau=0$ there,
then \eqref{eq:neck-geodesic} and \eqref{eq:variational-end-germ} give
$\ddot\tau=-\kappa p_1^2/2$.  If $p_1\ne0$, the boundary point is a
strict local maximum of $\tau$, and nearby times have $\tau<0$.
If $p_1=0$, then the scalar equation
$\ddot\tau=-F'(\tau)/2$ has the constant solution $\tau=0$ with these
initial data.  Smooth ODE uniqueness forces this solution locally.
Neither possibility can be an endpoint of a nonempty open-neck
component.  This proves transversality and also handles a boundary
grazing time.  In the second case the corresponding geodesic is the
fiber geodesic in the boundary torus, so uniqueness keeps it in that
torus for all time.

If $v$ vanishes at an interior time and $F'$ also vanishes there,
uniqueness again makes $t$ constant.  Then the solution stays at this
height for all time, so the passage is $\mathbb R$.
Otherwise such a zero is a strict extremum of $t$.  Reflecting the
interval coordinate if necessary makes it a maximum $t_*$ with
$c:=F'(t_*)>0$.  Concavity gives $F'(t)\ge c$ for $0\le t\le t_*$.
On each side of the turn, the trajectory moves toward $t=0$ and cannot
turn again, because $\dot v=-F'(t)/2\le-c/2$.
It reaches $0$ in finite time: after the turn, for example,
$t(s)\le t_* -c(s-s_*)^2/4$ as long as it stays in the neck.
Time reversal and uniqueness make the two halves equal.  Translating
the parameter gives \eqref{eq:turn-data}.  Concavity also gives
$F'(t)\le F'(0)=\kappa p_1^2$, proving
\eqref{eq:turn-concavity}.  Since $F'(t_*)>0$, $p_1$ cannot be zero.
Finally,
\[
 v_0^2=F(t_*)-F(0)=\int_0^{t_*}F'(u)\,du
          \le\kappa p_1^2t_*\le\kappa p_1^2,
\]
which proves \eqref{eq:turn-incidence}.

It remains to classify passages with $v$ nowhere zero.  Here $t$ is
strictly monotone.  If both endpoints are finite, transversality and
monotonicity force the two different boundary tori.  At an infinite
endpoint, bounded monotonicity gives a limit $t_\infty$.
The identity $v^2=1-F(t)$ shows that $|v|$ has a limit, which must be
zero because $t$ is bounded.  Thus $F(t_\infty)=1$.
Since $\dot v$ tends to $-F'(t_\infty)/2$, this limit must also be zero,
or $v$ could not tend to zero.  If $t_\infty$ were a boundary height,
\eqref{eq:variational-end-germ} and $F'(t_\infty)=0$ would force
$p_1=0$ in that end frame.  Then $F$ would be constant and equal to $1$
throughout its end portion, giving $v=0$ there, a contradiction.
Hence the limiting height is interior.

A nonconstant monotone passage cannot have two infinite endpoints.
Indeed, its two limiting heights would be distinct, with $F=1$ at
both.  Concavity would give $F\ge1$ between them, whereas
\eqref{eq:neck-geodesic} gives $F<1$ at every traversed height.
This contradiction finishes the classification.
\end{proof}

\subsection{Passages requiring no turning estimate}

\begin{lemma}\label{lem:elementary-passages}
Let $V$ be the restriction of a smooth compactly supported field along
$\gamma$ to an entire neck passage $J$.  If $J=\mathbb R$, or $v$ never
vanishes on $J$, or the passage is a turn contained in one end portion
$0\le\tau\le1/4$, then
\[
 I_J^\perp(V)\ge0.
\]
\end{lemma}

\begin{proof}
Suppose first that $J=\mathbb R$.  Replace $V$ by $V^\perp$.
This is a tangential change preserving compact support, and the new
tangential component $q_{V^\perp}$ is zero.  Apply
\eqref{eq:coordinate-index} on a compact interval whose endpoints lie
outside that support.  There is no boundary contribution, so
\[
 I_J^\perp(V)=I_J(V^\perp)=E_J(V^\perp)\ge0.
\]

Suppose next that $v$ is nowhere zero.  On the passage, replace $V$ by
\[
 \widetilde V=V-\frac{w}{v}T.
\]
Its interval component is identically zero.  This change is smooth
through every finite endpoint by transversality.  At an infinite
endpoint it is identically zero for sufficiently large $|s|$, since
the original field is zero there.  In particular, a possible limit
$v\to0$ at infinity causes no division problem: all nonzero values of
the numerator lie in a compact time interval where $v\ne0$.
The tangential change can, if needed, be extended through the finite
endpoints and cut off to a compactly supported change along the whole
geodesic.  No estimate uniform in $1/|v|$ is being asserted or used.

Write $\widetilde V=(0,\widetilde z)$.  Its tangential component is
$q_{\widetilde V}=p^{\mathsf T}\widetilde z$.  The endpoint signs in
Lemma~\ref{lem:coordinate-index}, with endpoints beyond the support
at each infinite end, give
\begin{align*}
 I_J^\perp(V)=I_J^\perp(\widetilde V)
 &\ge\int_J\bigl(|\dot{\widetilde z}|_Q^2
                    -|p^{\mathsf T}\dot{\widetilde z}|^2\bigr)\,ds\\
 &\ge\int_J(1-F(t))|\dot{\widetilde z}|_Q^2\,ds\ge0.
\end{align*}
The second inequality is Cauchy--Schwarz in the positive definite
matrix $Q$, using
$(p^{\mathsf T}\xi)^2\le(p^{\mathsf T}Pp)(\xi^{\mathsf T}Q\xi)$.

Finally, on an end portion sectional curvature is nonpositive.  The
normal-index integrand is therefore pointwise nonnegative, by
Lemma~\ref{lem:tangential-invariance}, which proves the last case.
\end{proof}

Only the finite turning excursions that leave an end portion remain.
We next normalize their fields without imposing any condition on the
original endpoint values.

\begin{lemma}[A gauge on a turning excursion]\label{lem:turning-gauge}
Consider a turning passage with data \eqref{eq:turn-data}--
\eqref{eq:turn-incidence}.  For every smooth field $V$ on $[0,L]$,
there is a tangentially equivalent smooth field, denoted again by
$V=(w,z)$, such that
\[
 w(0)=w(L)=0,\qquad \langle V,T\rangle\ \text{is affine in }s.
\]
Put $\Delta z=z(L)-z(0)$, $e=(1,0)^{\mathsf T}$, and
$\mu=\kappa v_0/4>0$.  For this field,
\begin{align}
 I_{[0,L]}^\perp(V)
 &=E_{[0,L]}(V)+\frac{\kappa v_0}{2}
       \bigl(z_1(0)^2+z_1(L)^2\bigr)
       -\frac{(p^{\mathsf T}\Delta z)^2}{L}
       \label{eq:turning-index-identity}\\
 &\ge E_{[0,L]}(V)+\mu(e^{\mathsf T}\Delta z)^2
       -\frac{(p^{\mathsf T}\Delta z)^2}{L}.
       \label{eq:turning-index-reduction}
\end{align}
The same left side is the normal index of the original field.
\end{lemma}

\begin{proof}
For the original field, let $q=\langle V,T\rangle$ and set
\[
 h_0=-\frac{w(0)}{v_0},\qquad h_L=\frac{w(L)}{v_0}.
\]
Let $\ell$ be the affine function with
$\ell(0)=q(0)+h_0$ and $\ell(L)=q(L)+h_L$.
The single change
\[
 \widetilde V=V+hT,\qquad h=\ell-q,
\]
has zero interval component at both endpoints and satisfies
$\langle\widetilde V,T\rangle=\ell$ everywhere.
This proves simultaneous solvability of the two requirements for
arbitrary original endpoint data.  If the original field is given on
the complete geodesic, multiplying $\ell-q$ by a smooth compactly
supported cutoff equal to $1$ on $[0,L]$ extends this to a compactly
supported tangential change; the identities on the excursion are
unchanged.

Drop the tilde.  Since the interval component is zero at the
endpoints, $q(0)=p^{\mathsf T}z(0)$ and
$q(L)=p^{\mathsf T}z(L)$, so its affine slope is
$p^{\mathsf T}\Delta z/L$.
At $s=0$ and $s=L$, \eqref{eq:boundary-christoffel} gives respectively
$\kappa v_0z_1(0)^2/2$ and $-\kappa v_0z_1(L)^2/2$ for
$\langle T,\Gamma(V,V)\rangle$.
Substitution in \eqref{eq:coordinate-index} and
\eqref{eq:normal-index-definition} proves
\eqref{eq:turning-index-identity}.  The inequality
$a^2+b^2\ge(b-a)^2/2$ then gives
\eqref{eq:turning-index-reduction}.
Tangential invariance returns both statements to the original
normal index form.
\end{proof}

The gauges in this section are used separately on each passage.
There is no assertion that gauges chosen on different passages match
at their endpoints.  Every inequality concerns the unchanged normal
index integral of the original field, which is the quantity to be
summed in the global argument.
